% ==============================================================================
% DECLARAÇÃO DE USO DE IA GENERATIVA
% Título sem indicativo numérico (centralizado), após o corpo textual e antes
% das Referências. Manter somente os itens efetivamente aplicáveis.
% ==============================================================================
\chapter*{Declaração de uso de IA generativa}
\phantomsection
\addcontentsline{toc}{chapter}{Declaração de uso de IA generativa}
\markboth{DECLARAÇÃO DE USO DE IA GENERATIVA}{}
\label{ch:declaracao_ia_generativa}
A declaração de uso de IA generativa tem como finalidade formalizar a integridade ética, o rigor metodológico e a transparência científica do trabalho, documentando a extensão da intervenção de ferramentas algorítmicas ao longo do processo de pesquisa e redação. O seu papel estrutural é discriminar com precisão os modelos e plataformas empregados (especificando versões, parâmetros ou finalidades de suporte), delimitar as tarefas assistidas (como revisão estilística, estruturação preliminar de código, tradução ou exploração conceitual) e atestar a primazia da autoria humana. Esse registro assegura que a responsabilidade crítica, analítica e epistemológica sobre as conclusões, a verificação de alucinações e a exatidão das referências permaneçam integralmente sob o domínio do autor, alinhando a dissertação às diretrizes de governança ética e integridade acadêmica estabelecidas pela comunidade científica e pelas agências de fomento.

Para cada intervenção mediada por modelos computacionais, reportam-se: a \textit{localização} (em que proposições ou seções a ferramenta introduziu, selecionou ou modificou conteúdo material), a \textit{extensão} (a natureza e o nível da operação realizada) e a \textit{verificação} (o procedimento pelo qual o autor examinou criticamente o resultado antes de assumir a responsabilidade pela sua incorporação).

\section*{Ferramentas utilizadas}
% Manter apenas os itens correspondentes às ferramentas efetivamente utilizadas.
% Repetir o modelo adequado quantas vezes forem necessárias.
\begin{itemize}
  % Modelo 1: execução local ou em rede local
  \item \textbf{\preencher{nome e versão do modelo}} (\preencher{empresa}), executado em infraestrutura computacional \preencher{local, na própria estação de trabalho, ou em servidor da rede local institucional} sob o ambiente \textit{\preencher{nome e versão do ambiente de execução, por exemplo, llama.cpp, Ollama ou LM Studio}}, sem transmissão de dados ou telemetria para servidores externos, durante sessões conduzidas em \preencher{mês e ano}.

        % Modelo 2: acesso por interface de chat hospedada pelo fornecedor
  \item \textbf{\preencher{nome e versão do modelo}} (\preencher{empresa}), acessado pela interface de chat \textit{\preencher{nome da plataforma, por exemplo, claude.ai, chatgpt.com ou gemini.google.com}}, em plano \preencher{gratuito, individual ou institucional}, com mecanismo de busca na web \preencher{ativado ou desativado}\preencher{, se ativado: para averiguação de registros externos} e com o uso dos dados para treinamento do modelo \preencher{habilitado ou desabilitado}, durante sessões conduzidas em \preencher{mês e ano}.
\end{itemize}

\section*{Intervenções e sua incidência sobre as afirmações}
As intervenções realizadas restringiram-se ao suporte linguístico e à verificação documental, não incidindo de forma material sobre as afirmações substantivas que fundamentam as conclusões do trabalho:
\begin{enumerate}
  \item \textbf{Sugestões de concisão e clareza textual} (\preencher{nome e versão}).
        \textit{Localização:} intervenções pontuais e dispersas ao longo \preencher{das seções ou capítulos, com a respectiva numeração}.
        \textit{Extensão:} reformulação estilística e ajuste sintático de orações preexistentes.
        \textit{Verificação:} as sugestões foram incorporadas pelo autor unicamente quando preservaram de modo estrito o escopo inferencial e a qualificação das premissas originais, tendo sido descartadas quaisquer formulações que alterassem a força das alegações.

  \item \textbf{Conferência de metadados bibliográficos e fontes} (\preencher{nome e versão}).
        \textit{Localização:} seção de Referências e confrontação de citações no texto.
        \textit{Extensão:} validação de identificadores persistentes (DOIs), volume, número, intervalos de paginação e correspondência com fontes primárias.
        \textit{Verificação:} cada entrada bibliográfica e metadado sugerido foram conferidos individualmente pelo autor, diretamente junto aos repositórios e às bases originais indexadas.
\end{enumerate}

\section*{Delimitação da autoria e produção não mediada por IA}
A formulação do problema de investigação, \preencher{tese, hipótese ou argumento central do trabalho}, \preencher{distinções teóricas, modelos conceituais ou taxonomias propostas}, \preencher{critérios, categorias analíticas ou procedimentos metodológicos definidos pelo autor} e \preencher{tabelas, figuras ou quadros conceituais de autoria própria, com a respectiva numeração} foram idealizados, estruturados e desenvolvidos de forma autônoma pelo autor, antes do emprego de quaisquer ferramentas computacionais.

A ferramenta \preencher{nome e versão} foi empregada na etapa de \preencher{busca, triagem ou seleção} da literatura, cabendo ao autor a verificação integral das fontes e a decisão final de inclusão.

\section*{Comandos textuais utilizados (\textit{Prompts})}
\begin{itemize}
  \item \textbf{Sugestões de concisão e clareza} (\preencher{nome e versão}):\\
        \textit{``Sugira opções de reformulação sintática para aumentar a concisão e a clareza do excerto a seguir, preservando integralmente o alcance das afirmações e sem introduzir novas proposições: [texto].''}

  \item \textbf{Conferência de metadados bibliográficos e fontes} (\preencher{nome e versão}):\\
        \textit{``Examine a lista de referências anexa e valide a exatidão dos respectivos metadados bibliográficos (especificamente autoria, título, periódico, ano, volume, número, páginas e DOI) e sua correspondência com os registros publicados: [lista de referências].''}
\end{itemize}
